\section{Divisorial witnesses for flipped surfaces}
\label{sec:witnesses}

A surface in the flipped locus supplies a divisor whose discrepancy has
just increased.  The fixed rational nef data constrain the new discrepancy
to a finite set.  The construction requires smoothness at the general
point of the flipped surface.  Generalized canonicity alone does not
supply that smoothness; strict increase on the target exceptional locus
does.  We first establish this local consequence and then construct a
global divisor by blowing up the surface.
For ordinary canonical pairs, this use of terminality near a flipped
surface appears in Fujino \cite[Lemma~2.2]{Fujino2004}.
The proof here includes the nef-trace defect and checks the local
singularity condition from the hypotheses on global places.

\begin{proposition}[Terminality near a flipped surface]
\label{prop:ordinary-terminal}
Let $T\dashrightarrow T^+$ be a step in the sequence of
Theorem~\ref{thm:termination}, and let $S\subset L^+$ be an irreducible
flipped surface.  There is an analytic open subset $U\subset T^+$
meeting $S$ on which the underlying space has ordinary terminal
singularities.  Consequently $T^+$ is smooth at a general point of $S$.
\end{proposition}

\begin{proof}
Choose a smooth model $r\colon W\to T^+$, projective over $T^+$ and
carrying the fixed nef data, with divisorial exceptional locus as in
Proposition~\ref{prop:common-models}.  Using the divisors of
\eqref{eq:boundary-and-defect}, write the ordinary relative canonical
equation as
\begin{equation}
 \Theta_W:=r^*K_{T^+}-K_W
   =\Delta_W-r^*B_{T^+}-J_W.
 \label{eq:ordinary-boundary}
\end{equation}
The pullback $r^*B_{T^+}$ is effective: after clearing a Cartier
multiple, a local equation for the effective boundary is holomorphic
by normality, and remains holomorphic after pullback.
Lemma~\ref{lem:nef-defect} gives $J_W\geq0$.  Generalized canonicity
therefore makes every exceptional coefficient of $\Theta_W$
nonpositive.  Its nonexceptional coefficients are zero, since it is
an ordinary relative canonical divisor.

An exceptional prime $P$ on $W$ whose image contains $S$ has image
exactly $S$: its irreducible image has dimension at most two.
Proposition~\ref{prop:comparison} gives
\[
 a(P;T^+,B_{T^+}+\bM)>a(P;T,B_T+\bM)\geq1,
\]
because its target centre lies in $L^+$.  Thus its coefficient in
$\Delta_W$, and hence in $\Theta_W$, is strictly negative.

There are only finitely many exceptional prime divisors on $W$.
Delete from $T^+$ the images of those with coefficient zero in
$\Theta_W$, and call the resulting open set $U$.  Proper mapping
makes the deleted set closed analytic.  None of these images contains
$S$, so their finite union cannot contain the irreducible surface $S$;
in particular, $U\cap S$ is a dense open subset of $S$.

We check ordinary terminality on $U$ using local analytic places.
Restrict the resolution and \eqref{eq:ordinary-boundary} over $U$.
Every local component of its exceptional divisor now has strictly
negative coefficient in $\Theta_W$.  If a local exceptional place
$E$ over $U$ has divisorial centre on this resolution, it is the place
of that divisor.  The divisor must be exceptional over $U$, since
otherwise $E$ would not be exceptional.  Its negative coefficient
therefore gives ordinary log discrepancy greater than one.

For any remaining local exceptional place, take a common local
resolution.  Its centre on the smooth space $r^{-1}(U)$ has
codimension at least two.  Since $\Theta_W\leq0$,
Lemma~\ref{lem:jacobian} gives
\[
 a(E;U,0)=a(E;r^{-1}(U),\Theta_W|_{r^{-1}(U)})
       \geq a(E;r^{-1}(U),0)\geq2.
\]
This checks all local exceptional places and proves ordinary
terminality.  The canonical rational line bundle restricts to $U$,
so Lemma~\ref{lem:terminal-smooth} applies there.  Its singular locus
has codimension at least three and cannot contain the surface
$S\cap U$.  Thus $T^+$ is smooth at a general point of $S$.
\end{proof}

\subsection{The blowup divisor and its discrepancy}

Fix an integer $m>0$ such that $mM'$ is a Cartier divisor and $mb$ is an
integer for every positive coefficient $b$ of $B_0$.  The same integer
clears all boundary coefficients on every model, since the boundaries are
strict transforms.  Moreover
\begin{equation}
 mM_T\text{ is an integral Weil divisor on every model }T.
 \label{eq:integral-nef-traces}
\end{equation}
Indeed, on a common model carrying the data, $mM_W$ is the pullback of the
integral Cartier divisor $mM'$, and pushforward preserves integral Weil
coefficients.  Although \eqref{eq:integral-nef-traces} does not assert
global Cartierness of $mM_T$, it implies Cartierness on the smooth locus
of $T$.

\begin{proposition}[A witness for every flipped surface]
\label{prop:witnesses}
Consider the $i$-th step of a sequence as in Theorem~\ref{thm:termination}, with target
$T^+=X_{i+1}$.  Let $S$ be an irreducible surface contained in its flipped
locus.  There is a global prime place $E_S$, with centre $S$ on $T^+$,
obtained by blowing up $S$ at its general smooth point.  It is exceptional
over every model of the sequence.  Write
\[
 B_{i+1}=\sum_{\ell} b_{\ell}H_{\ell},
 \qquad b_{\ell}>0,
\]
and let $r\colon W\to T^+$ be a smooth model that is projective over
$T^+$ and carries both $E_S$ and the fixed nef data.  With
$J_W=r^*M_{i+1}-M_W$, one has
\begin{equation}
 a(E_S;X_{i+1},B_{i+1}+\bM)
 =2-\sum_{\ell}b_{\ell}\operatorname{mult}_S(H_{\ell})
       -\operatorname{coeff}_{E_S}(J_W).
 \label{eq:witness-discrepancy}
\end{equation}
In particular,
\begin{equation}
 a(E_S;X_{i+1},B_{i+1}+\bM)
 \in\Lambda_m:=(1,2]\cap\tfrac1m\Z
 =\left\{1+\tfrac1m,1+\tfrac2m,\ldots,2\right\}.
 \label{eq:witness-grid}
\end{equation}
Its discrepancy on the source is strictly below two.  If $S$ is contained
in a component of $B_{i+1}$ of coefficient $b>0$, then more precisely
\begin{equation}
 a(E_S;X_i,B_i+\bM)
 < a(E_S;X_{i+1},B_{i+1}+\bM)\leq2-b.
 \label{eq:witness-threshold}
\end{equation}
\end{proposition}

\begin{proof}
By Proposition~\ref{prop:ordinary-terminal}, a general point of the
surface $S$ is a smooth point of both $T^+$ and $S$.
Blow up the global coherent
ideal of the reduced analytic subspace $S$.  Over the open subset where
$T^+$ and $S$ are smooth, this is the blowup of a smooth codimension-two
centre, with one exceptional prime dominating that open subset of $S$.
After normalization the preimage of $S$ is a compact analytic subset
with finitely many irreducible components.  One of these contains this
exceptional divisor and dominates $S$, by proper mapping; the smooth
calculation gives uniqueness.  It is a global prime divisor, and we
denote its place by $E_S$.  Resolve further and use
Proposition~\ref{prop:common-models} to obtain a smooth model $W$,
projective over $T^+$, carrying it and the nef data.  The centre on $T^+$ is $S$, so it is
exceptional there.  Smallness of every step preserves exceptionalness,
as in Proposition~\ref{prop:comparison}, and therefore it is exceptional
over every $X_j$.

The ordinary log discrepancy of the blowup at a smooth codimension-two
centre is two.  Pulling back a boundary component $H_{\ell}$ subtracts
$b_{\ell}\operatorname{mult}_S(H_{\ell})$, where the generic
multiplicity is a nonnegative integer and is zero unless $S\subset
H_{\ell}$.  The identity
\[
 \Delta_W=r^*(K_{T^+}+B_{i+1})-K_W+J_W
\]
then proves \eqref{eq:witness-discrepancy}.  This computation only uses
the smooth general point of the centre.  The boundary components may be
singular or intersect one another along $S$.

Lemma~\ref{lem:nef-defect} gives
$\operatorname{coeff}_{E_S}(J_W)\geq0$.  It also lies in
$\tfrac1m\Z$.  To verify this denominator claim, restrict to a smooth
open subset of $T^+$ meeting a general point of $S$.
By \eqref{eq:integral-nef-traces}, $mM_{i+1}$ is an integral Cartier
divisor there, so its pullback has integral order along $E_S$.
The divisor $mM_W$ is integral Cartier as the pullback of $mM'$.
Consequently the coefficient of
\[
 mJ_W=r^*(mM_{i+1})-mM_W
\]
along $E_S$ is an integer.  This is a local calculation near the general
point of $S$, and does not require $mM_{i+1}$ to be Cartier elsewhere.

Every term subtracted from two in \eqref{eq:witness-discrepancy} is
nonnegative and belongs to $\tfrac1m\Z$.  Since the target centre is
contained in $L_i^+$, strict increase in
Proposition~\ref{prop:comparison} and generalized canonicity on the
source give
\[
 a(E_S;X_{i+1},B_{i+1}+\bM)
 >a(E_S;X_i,B_i+\bM)\geq1.
\]
These facts prove \eqref{eq:witness-grid}.  If $S$ is contained in a component
of coefficient $b$, its generic multiplicity is at least one, so
\eqref{eq:witness-discrepancy} bounds the target discrepancy by $2-b$.
This proves \eqref{eq:witness-threshold}; without that containment the
same argument gives the strict source bound below two.
\end{proof}

Call a place chosen by Proposition~\ref{prop:witnesses} a \emph{witness}
for the corresponding flipped surface.  The possible target discrepancies
of witnesses form the fixed finite set $\Lambda_m$.  A place that is used
again must reach a strictly larger member of this set.

\begin{lemma}[Finite use of each witness]
\label{lem:finite-use}
Let $I$ be any set of step indices in a sequence as in Theorem~\ref{thm:termination}.  For every
$i\in I$, choose a flipped surface of the $i$-th step and a witness $E_i$
as in Proposition~\ref{prop:witnesses}.  Then each fixed place $E$
occurs among the $E_i$ at most $m$ times.  Consequently, if all the
$E_i$ belong to a fixed finite collection $\mathcal F$, then
$\#I\leq m\,\#\mathcal F$.
\end{lemma}

\begin{proof}
Suppose $i<j$ are two indices for which $E_i=E_j=E$.  Monotonicity and
the strict increase when $E$ is used at step $j$ give
\[
 a(E;X_{i+1},B_{i+1}+\bM)
 \leq a(E;X_j,B_j+\bM)
 <a(E;X_{j+1},B_{j+1}+\bM).
\]
Thus the target discrepancies at all uses of this fixed place are
strictly increasing elements of $\Lambda_m$.  This set has $m$ elements,
so there can be at most $m$ uses.  Summing this bound over
$\mathcal F$ proves the last assertion.  No discreteness of the
discrepancy at intervening steps is required.
\end{proof}
